% ============================================================================
%  CHƯƠNG 0: TÓM TẮT
% ============================================================================

\section*{Tóm tắt}
\addcontentsline{toc}{section}{Tóm tắt}

Đồ án xây dựng hệ thống trích xuất tên cửa hàng, ngày giao dịch, tổng tiền và địa chỉ từ ảnh biên lai tiếng Việt. Ba phương pháp được triển khai: Baseline dùng OCR kết hợp luật, \model{Donut} sinh chuỗi trực tiếp từ ảnh, và \model{LayoutXLM} phân loại token dựa trên văn bản, bố cục và ảnh. Bộ dữ liệu processed có 1.559 mẫu; Donut dùng được toàn bộ 1.091 mẫu train, còn LayoutXLM dùng 784 mẫu train có hộp bao. Các prediction artifact đều bao phủ cùng 234 ID test.

Để tránh điểm số bị ảnh hưởng bởi nhãn rỗng, báo cáo đưa ra cả kết quả trên toàn bộ mẫu lẫn kết quả present-only (chỉ xét mẫu có ground-truth khác rỗng của từng trường). Trong ba phương pháp so sánh, \model{LayoutXLM} cho kết quả tốt nhất với Macro EM 42,96\% và Macro NES 63,89\% trên present-only; Baseline đạt 40,41\% và 55,47\%; còn \model{Donut} được suy luận với giới hạn sinh 768 token nhưng vẫn cho hiệu năng rất thấp (Macro EM 0,00\%) vì hiện tượng lặp từ và model collapse khi huấn luyện end-to-end trên tập dữ liệu nhỏ. Thí nghiệm Oracle OCR cho thấy nếu có OCR hoàn hảo, hiệu năng LayoutXLM tăng lên Macro EM 66,97\% trên present-only --- điều này chỉ ra OCR chính là nút thắt cổ chai lớn nhất của pipeline. Báo cáo cũng cập nhật đầy đủ log huấn luyện, phân tích độ nhạy hệ tọa độ và hoàn thiện bằng chứng thực nghiệm.

\vspace{1em}
\noindent\textbf{Từ khoá:} trích xuất thông tin, biên lai, Donut, LayoutXLM, OCR, tiếng Việt.
